% =============================================================================
% Denis Stepanov — CV (Russian)
% Layout: sidebar on the LEFT (education, achievements, publications,
% languages); main column on the RIGHT (experiences with bullets).
% Built on AltaCV (LianTze Lim, v1.1.5). Compile with: tectonic resume_ru.tex
% =============================================================================

\documentclass[10pt,a4paper,ragged2e]{altacv}

% ─── Page layout ──────────────────────────────────────────────────────────────
% Sidebar lives in the LEFT margin (left=10cm reserves space for it).
\geometry{
  left=10cm, right=2cm,
  marginparwidth=6.8cm, marginparsep=1.2cm,
  top=0.6cm, bottom=1.25cm
}

% Put \marginpar content on the LEFT instead of the default right.
\reversemarginpar

% ─── Packages ─────────────────────────────────────────────────────────────────
\usepackage[super]{nth}
\usepackage{float}
\usepackage{changepage}  % already loaded by altacv but re-declaring for clarity
% Note: altacv.cls has been locally patched to load biblatex with backend=bibtex
% (this CV doesn't print a bibliography, so biber doesn't need to run).

% Redefine \fullwidth so it extends LEFT into the sidebar gutter instead of
% right (matches the mirrored layout). AltaCV's original definition extends
% the right side; we extend the left.
\renewenvironment{fullwidth}{%
  \begin{adjustwidth}{\dimexpr-\marginparwidth-\marginparsep\relax}{}}%
  {\end{adjustwidth}}

% ─── Fonts & language ─────────────────────────────────────────────────────────
% Tectonic uses XeTeX, so the xetex branch is the active one.
\ifxetexorluatex
  \usepackage{fontspec}
  \usepackage[english,russian]{babel}
  \setmainfont{Carlito}[Script=Cyrillic]
\else
  \usepackage[utf8]{inputenc}
  \usepackage[T2A,T1]{fontenc}
  \usepackage[russian]{babel}
  \usepackage[default]{lato}
\fi

% ─── Theme colors ─────────────────────────────────────────────────────────────
\definecolor{VividPurple}{HTML}{195e84}  % accent / headings
\definecolor{SlateGrey}{HTML}{2E2E2E}    % body
\colorlet{heading}{VividPurple}
\colorlet{accent}{VividPurple}
\colorlet{emphasis}{SlateGrey}
\colorlet{body}{SlateGrey}

% ─── Bullets ──────────────────────────────────────────────────────────────────
\renewcommand{\itemmarker}{{\small\textbullet}}
\renewcommand{\ratingmarker}{\faCircle}

% Slightly smaller font inside \begin{itemize}.
\AtBeginEnvironment{itemize}{\small}

% Custom contact field for Telegram (AltaCV has no built-in one).
\newcommand{\telegram}[1]{\printinfo{\faPaperPlane}{#1}}


\begin{document}

% =============================================================================
% HEADER
% =============================================================================
\name{Денис Степанов}
\tagline{Mobile \& Frontend Developer}
\photo{3.3cm}{profile.jpg}
\personalinfo{
  \email{dentep78@gmail.com}
  \phone{+79933518359}
  \telegram{t.me/denistepp}
  \location{Дананг, Вьетнам (GMT+7)}
  \homepage{dentep.github.io/}
  \github{https://github.com/dentep}
}

\begin{fullwidth}
  \makecvheader
\end{fullwidth}

\vspace{0.4cm}
\begin{fullwidth}
\small Разработчик мобильных приложений с 5+ годами опыта в React Native.
Веду проекты от прототипа до релиза в App Store и Google Play.
Открыт к удалённой работе и релокации.
\end{fullwidth}
\vspace{0.2cm}


% =============================================================================
% MAIN COLUMN (right) — Опыт работы
% Sidebar content lives in page1sidebar.tex (Образование, Достижения,
% Публикации, Языки), attached to the first \cvsection below.
% =============================================================================

\cvsection[page1sidebar]{Опыт работы (7+ лет)}

% ─── MD Audit ────────────────────────────────────────────────────────────────
\cvevent
  {Frontend Разработчик (Vue)}
  {MD Audit}
  {Май 2025 -- Сейчас}
  {Удалённо}
\begin{itemize}
  \item Разработка веб-интерфейсов на Vue для внутренних аудиторских сервисов.
\end{itemize}
\divider

% ─── Softline AG ─────────────────────────────────────────────────────────────
\cvevent
  {Mobile Dev Lead}
  {Softline AG}
  {Июнь 2022 -- Май 2025}
  {Удалённо}
\begin{itemize}
  \item Возглавлял команду мобильной разработки, выпустил 100+ релизов в production.
  \item Снизил частоту крашей и улучшил стабильность приложений на 94\% за 3 года.
  \item Внедрил процессы код-ревью, CI/CD и автотестирования.
  \item Участвовал в найме и менторинге младших разработчиков.
\end{itemize}
\divider

% ─── Briskly ─────────────────────────────────────────────────────────────────
\cvevent
  {Тимлид Разработки Мобильных Приложений}
  {Briskly}
  {Март 2021 -- Июль 2022}
  {Москва, Россия}
\begin{itemize}
  \item Руководил разработкой B2C приложения B-PAY (300 000+ загрузок).
  \item Выпустил CRM-приложение Briskly Business для продавцов розничной сети.
  \item Координировал работу команды из 4 разработчиков.
\end{itemize}
\divider

% ─── Sycret ──────────────────────────────────────────────────────────────────
\cvevent
  {Senior Разработчик Мобильных Приложений}
  {Sycret}
  {Май 2020 -- Апрель 2021}
  {Москва, Россия}
\begin{itemize}
  \item Мигрировал CRM-приложение с C\#/Xamarin на React Native.
  \item Поддерживал линейку из 22 white-label приложений на iOS и Android.
  \item Реализовал интеграции с REST API, push-уведомлениями и аналитикой.
\end{itemize}
\divider

% ─── UNNC ────────────────────────────────────────────────────────────────────
\cvevent
  {Дизайнер Виртуальной Среды (VE)}
  {University of Nottingham Ningbo China}
  {Июнь 2019 -- Май 2020}
  {Нинбо, Китай}
\begin{itemize}
  \item Создал 4D VR-окружение для обучения студентов гражданскому строительству.
  \item Результаты работы опубликованы в материалах конференций IEEE.
\end{itemize}

\end{document}
